%%%%%%%%%%%%%%%%%%%%%%foreword.tex%%%%%%%%%%%%%%%%%%%%%%%%%%%%%%%%%
% sample foreword
%
% Use this file as a template for your own input.
%
%%%%%%%%%%%%%%%%%%%%%%%% Springer %%%%%%%%%%%%%%%%%%%%%%%%%%

%\foreword

\chapter*{Prólogo}
\markboth{Prólogo}{Prólogo}
\addcontentsline{toc}{chapter}{Prólogo}

El estudio de los autómatas celulares es una de las áreas más interesantes dentro de la teoría de la computación y de los sistemas dinámicos discretos. A pesar de que parte de reglas relativamente simples, los modelos han demostrado ser capaces de producir comportamientos complejos, llegando a representar procesos que se relacionan con computación universal. Debido a esto, los autómatas celulares continúan siendo objeto de investigación en diferentes ramas de matemáticas y la computación.

El presente Trabajo Terminal toca temas justamente en esa línea de investigación, enfocándose en el análisis de la computabilidad en autómatas celulares. A lo largo del presente documento, los autores desarrollan una aproximación fundamentada en teoría sobre topología discreta, dinámica de sistemas y teoría de autómatas, lo cual permite interpretar fenómenos emergentes dentro de espacios celulares bidimensionales.
Uno de los aspectos a destacar en este trabajo es el esfuerzo por mantener una estructura rigurosa en la exposición de las ideas. El documento no se limita únicamente a describir comportamientos que se observaron en la experimentación, sino que intenta justificar las propiedades computacionales asociadas a reglas complejas como \textit{Spiral} y \textit{Beehive}. Este enfoque requiere una comprensión sólida de los fundamentos teóricos involucrados.
Un punto importante es la realización establecida entre las colisiones de estructuras y la posibilidad de sintetizar operaciones lógicas. La idea de usar gliders y dinámicas emergentes como mecanismos de procesamiento de información representa una aproximación distinta al paradigma clásico de Von Neumann. La dirección general de la investigación se encuentra claramente definida.
También es evidente el esfuerzo que se puso para mantener la coherencia entre la parte teórica y la implementación propuesta. Validar los modelos mediante simulaciones aporta un componente práctico relevante que permite trasladar conceptos teóricos a un entorno “tangible”.
El trabajo aquí presentado constituye una aportación valiosa al campo de estudio al que se enfoca, que puede servir como base para futuras investigaciones relacionadas con autómatas celulares y modelos alternativos de computación.


\vspace{\baselineskip}
\begin{flushright}\noindent
Mayo de 2026\hfill {\it Juan Alberto Rugerio Arceo}\\
\end{flushright}


